
% --------------------------------------------------------------------------
% Acronym and glossary categories
%
% Add, remove, or rename categories here. The corresponding sections in
% the List of Acronyms are generated automatically.
% --------------------------------------------------------------------------

\DeclareAcronymType
    {phystype}{phl}{pho}{phg}
    {Physical Quantities and Concepts}

\DeclareAcronymType
    {stattype}{stl}{sto}{stg}
    {Statistics and Mathematics}

\DeclareAcronymType
    {compmethodtype}{cml}{cmo}{cmg}
    {Computation and Analysis}

\DeclareAcronymType
    {orgtype}{orl}{oro}{org}
    {Scientific Organizations and Facilities}

\DeclareAcronymType
    {othertype}{otl}{oto}{otg}
    {Miscellaneous}
    
%%%%%%%%%%%%%%%%%%%%%%%%%%%%%%%%%%%%%%%%%%%%%%%%%%%%%%%%%%%%
% Definition of a macro for Glossary / Abbrev entries 
%  Some entries need cross linking between both sections
%  This macro will achieve this easily
%%%%%%%%%%%%%%%%%%%%%%%%%%%%%%%%%%%%%%%%%%%%%%%%%%%%%%%%%%%%
% Arguments:
% Paired abbreviation + glossary entry
% Optional args:
%   #1  options for \newabbreviation (type=..., etc.)
%   #2  abbrev label
%   #3  short form (singular)
%   #4  long form (singular, plain)
%   #5  glossary label
%   #6  glossary description
%   #7  [optional] abbreviation plural short (default: #3s)
%   #8  [optional] glossary plural (default: #4s)

\newcommand{\glsnouse}[1]{%
  \glslink{#1}{\glsxtrshort{#1}}%
}
\newcommand{\glsnousepl}[1]{%
  \glslink{#1}{\glsxtrshortpl{#1}}%
}
\newcommand{\glslinkpl}[1]{%
  \glslink{#1}{\glsentryplural{#1}}%
}

\makeatletter
\newcommand{\newabbrglspair}[8][]{%
  % --- defaults for plurals ---
  \def\abbrplural{#7}%
  \ifx\abbrplural\@empty
    \edef\abbrplural{#3s}%
  \fi
  \def\glsplural{#8}%
  \ifx\glsplural\@empty
    \edef\glsplural{#4s}%
  \fi

  % --- Abbreviation definition ---
  \newabbreviation[
    #1,
    shortplural=\abbrplural,
    longplural=\glsplural,
    user1={(see \glslink{#5}{Glossary})}
  ]{#2}{#3}{#4}%

  % --- Glossary entry definition ---
  \newglossaryentry{#5}{%
    name={#4},
    plural={\glsplural},
    description={#6 (Abbr.: \glsnouse{#2})}%
  }%
}
\makeatother

%%%%%%%%%%%%%%%%%%%%%%%%%%%%%%%%%%%%%%%%%%%%%%%%%%%%%%%%%%%%
% Double Abbreviation / Acronym / Initialisation AND glossary 
%  entries start here - note the 'type' will determine the 
%  section entry in the abbreviation list
%%%%%%%%%%%%%%%%%%%%%%%%%%%%%%%%%%%%%%%%%%%%%%%%%%%%%%%%%%%%

%%%%%%%%%%%%%%%%%%%%%%%%%%%%%%%%%%%%%%%%%%%%%%%%%%%%%%%%%%%%
% Paired abbreviation + glossary entries
%%%%%%%%%%%%%%%%%%%%%%%%%%%%%%%%%%%%%%%%%%%%%%%%%%%%%%%%%%%%

\newabbrglspair[type=phystype]
    {em}{EM}{electromagnetic}{electromagnetic-radiation}{%
    Radiation consisting of coupled electric and magnetic fields that
    propagate through space. Electromagnetic radiation spans a broad range
    of wavelengths and frequencies, including radio waves, visible light,
    ultraviolet radiation, X-rays, and gamma rays.
}{}{}

\newabbrglspair[type=phystype]
    {ke}{KE}{kinetic energy}{kinetic-energy}{%
    The energy associated with the motion of an object. In classical
    mechanics, the kinetic energy of an object of mass $m$ moving with
    speed $v$ is $K = \frac{1}{2}mv^2$.
}{}{}

\newabbrglspair[type=phystype]
    {pe}{PE}{potential energy}{potential-energy}{%
    Energy associated with the position or configuration of a physical
    system. Examples include gravitational potential energy and elastic
    potential energy.
}{}{}

\newabbrglspair[type=stattype]
    {pdf}{PDF}{probability density function}{probability-density-function}{%
    A function describing the relative likelihood that a continuous random
    variable takes a particular value. The integral of a probability
    density function over an interval gives the probability that the
    variable lies within that interval.
}{}{}

\newabbrglspair[type=stattype]
    {sd}{SD}{standard deviation}{standard-deviation}{%
    A measure of the dispersion of a set of values about their mean.
    For a probability distribution, the standard deviation is the square
    root of the variance.
}{}{}

\newabbrglspair[type=compmethodtype]
    {mc}{MC}{Monte Carlo}{monte-carlo-method}{%
    A class of computational methods that use repeated random sampling to
    estimate numerical quantities, propagate uncertainty, or explore
    complex parameter spaces.
}{}{}

%%%%%%%%%%%%%%%%%%%%%%%%%%%%%%%%%%%%%%%%%%%%%%%%%%%%%%%%%%%%
% Single Abbreviation / Acronym / Initialisation  
% i.e. no companion glossary entries start here 
%   - note the 'type' will determine the section entry
%%%%%%%%%%%%%%%%%%%%%%%%%%%%%%%%%%%%%%%%%%%%%%%%%%%%%%%%%%%%
%%%%%%%%%%%%%%%%%%%%%%%%%%%%%%%%%%%%%%%%%%%%%%%%%%%%%%%%%%%%
% Single abbreviations / acronyms
%%%%%%%%%%%%%%%%%%%%%%%%%%%%%%%%%%%%%%%%%%%%%%%%%%%%%%%%%%%%

% Physical quantities and concepts
\newabbreviation[type=phystype]{ac}{AC}{alternating current}
\newabbreviation[type=phystype]{dc}{DC}{direct current}
\newabbreviation[type=phystype]{emf}{EMF}{electromotive force}
\newabbreviation[type=phystype]{ir}{IR}{infrared}
\newabbreviation[type=phystype]{uv}{UV}{ultraviolet}

% Statistical and mathematical terms
\newabbreviation[type=stattype]{ci}{CI}{confidence interval}
\newabbreviation[type=stattype]{cdf}{CDF}{cumulative distribution function}
\newabbreviation[type=stattype]{mle}{MLE}{maximum likelihood estimation}
\newabbreviation[type=stattype]{pca}{PCA}{principal component analysis}
\newabbreviation[type=stattype]{rms}{RMS}{root mean square}

% Computing and analysis
\newabbreviation[type=compmethodtype]{api}{API}{application programming interface}
\newabbreviation[type=compmethodtype]{cpu}{CPU}{central processing unit}
\newabbreviation[type=compmethodtype]{gpu}{GPU}{graphics processing unit}
\newabbreviation[type=compmethodtype]{hpc}{HPC}{high-performance computing}
\newabbreviation[type=compmethodtype]{mcmc}{MCMC}{Markov chain Monte Carlo}

% Scientific organizations / facilities
\newabbreviation[type=orgtype]{aps}{APS}{American Physical Society}
\newabbreviation[type=orgtype]{nsf}{NSF}{National Science Foundation}
\newabbreviation[type=orgtype]{nist}{NIST}{National Institute of Standards and Technology}
\newabbreviation[type=orgtype]{utrgv}{UTRGV}{University of Texas Rio Grande Valley}

% Miscellaneous scientific usage
\newabbreviation[type=othertype]{doi}{DOI}{digital object identifier}
\newabbreviation[type=othertype]{si}{SI}{International System of Units}
\newabbreviation[type=othertype]{snr}{SNR}{signal-to-noise ratio}

%%%%%%%%%%%%%%%%%%%%%%%%%%%%%%%%%%%%%%%%%%%%%%%%%%%%%%%%%%%%
% Single Glossary entries (i.e. no companion abbreviation 
%   - start here
%%%%%%%%%%%%%%%%%%%%%%%%%%%%%%%%%%%%%%%%%%%%%%%%%%%%%%%%%%%%
%%%%%%%%%%%%%%%%%%%%%%%%%%%%%%%%%%%%%%%%%%%%%%%%%%%%%%%%%%%%
% Single glossary entries
%%%%%%%%%%%%%%%%%%%%%%%%%%%%%%%%%%%%%%%%%%%%%%%%%%%%%%%%%%%%

\newglossaryentry{conservation-energy}{%
    name={conservation of energy},
    description={%
    The principle that the total energy of an isolated system remains
    constant, although energy may be transferred between components or
    transformed from one form into another.
    }%
}

\newglossaryentry{momentum}{%
    name={momentum},
    description={%
    A vector quantity defined classically as the product of an object's
    mass and velocity, $\mathbf{p}=m\mathbf{v}$. For an isolated system,
    total momentum is conserved.
    }%
}

\newglossaryentry{uncertainty}{%
    name={measurement uncertainty},
    description={%
    A quantitative characterization of the range of values that may
    reasonably be attributed to a measured quantity. Measurement
    uncertainty may contain both random and systematic contributions.
    }%
}

\newglossaryentry{systematic-error}{%
    name={systematic error},
    description={%
    A persistent bias or offset in a measurement or analysis that does not
    average to zero when observations are repeated. Systematic effects may
    arise from calibration, instrumentation, methodology, or model
    assumptions. See also \glshyperlink{random-error}.
    }%
}

\newglossaryentry{random-error}{%
    name={random error},
    description={%
    Variation in repeated measurements arising from stochastic fluctuations
    or unpredictable influences. Random errors can often be reduced by
    increasing the number of independent measurements. See also
    \glshyperlink{systematic-error}.
    }%
}

\newglossaryentry{calibration}{%
    name={calibration},
    description={%
    The process of establishing the relationship between the output of an
    instrument or measurement system and known reference values.
    }%
}

\newglossaryentry{model}{%
    name={scientific model},
    description={%
    A mathematical, computational, or conceptual representation of a
    physical system used to describe observations, test hypotheses, or
    generate predictions.
    }%
}

\newglossaryentry{parameter}{%
    name={model parameter},
    description={%
    A quantity appearing in a \glshyperlink{model} whose value controls
    some aspect of the model's behavior and may be specified, measured, or
    inferred from data.
    }%
}

\newglossaryentry{residual}{%
    name={residual},
    description={%
    The difference between an observed or measured value and the
    corresponding value predicted by a model.
    }%
}

\newglossaryentry{normal-distribution}{%
    name={normal distribution},
    description={%
    A continuous probability distribution characterized by a mean and
    variance and having a symmetric, bell-shaped probability density
    function. It is also commonly known as the Gaussian distribution.
    }%
}

\newglossaryentry{correlation}{%
    name={correlation},
    description={%
    A statistical association between two variables. Correlation describes
    the degree to which the variables vary together but does not, by itself,
    establish a causal relationship.
    }%
}

\newglossaryentry{reproducibility}{%
    name={reproducibility},
    description={%
    The extent to which an analysis or result can be independently obtained
    using the same data, methods, and computational procedures.
    }%
}
